\chapter{Monitor ao vivo}

\eyebrow{sobreposição de gordura em tempo real para enquadramento}

\vspace{4pt}
O \textbf{Monitor ao vivo} exibe uma sobreposição de gordura em tempo real, de uma câmera ou de um
arquivo de vídeo, para auxiliar no enquadramento da carcaça durante a captura. Trata-se de um
\emph{guia de enquadramento}, não da medição oficial.

\section{Como usar}

\shot{live-monitor.jpeg}{O Monitor ao vivo, parado. Um seletor de origem (Câmera ao vivo /
Arquivo de vídeo) escolhe a fonte do quadro; \tele{Start} inicia a sobreposição e
\tele{Capture background} registra a cena vazia para a subtração. O estado do motor
e o dispositivo são exibidos no topo.}{live-monitor}

\begin{enumerate}[leftmargin=1.4em,itemsep=3pt]
  \item A origem é selecionada: \textbf{Câmera ao vivo} ou \textbf{Arquivo de vídeo}.
  \item O controle \tele{Start} inicia a reprodução do vídeo, e a sobreposição de gordura
        (ciano) o acompanha, com uma leitura ao vivo de FPS e de percentual de gordura no canto.
  \item Para uma câmera fixa, o controle \tele{Capture background} registra a cena vazia;
        o aplicativo a subtrai de cada quadro para que a sobreposição fique restrita à
        carcaça. A carcaça é então posicionada.
\end{enumerate}

\section{Como funciona e seus limites}

O monitor usa subtração de fundo clássica (rápida, $\sim$10\,ms) somada ao
modelo de gordura em resolução reduzida, atingindo taxas próximas do tempo real no
hardware testado. Ele \emph{não} usa a remoção de fundo precisa (porém lenta)
do pipeline de Análise.

\begin{caution}[title={Guia de enquadramento, não o número validado}]
A sobreposição ao vivo é aproximada e serve como auxílio visual para o enquadramento. O
número de gordura validado provém da aba de Análise, executada sobre uma imagem capturada com
remoção de fundo precisa.
\end{caution}
